% !PRINTED 321
\weberpage{321}
% [page alignment corrected 2026-08-07, 4x zoom strip 4/4 of idx 345] Three blocks that opened
% this fragment -- the tail of theorem XIII, the sentence "Wenn wir dies Theorem wiederholt
% anwenden…", and the whole of theorem XIV -- are PRINTED ON PAGE 320 and have been moved to
% i0345 where they belong. Page 321's top band begins, directly under its running head, with
% "Diese Sätze werden wesentlich verallgemeinert…".
%
% The %CONT>>/%<<CONT pair was internally consistent, so the split LOOKED sound; it was simply
% made at the wrong place on the page. Nothing mechanical can see this -- both fragments carried
% the right printed number and the volume compiled -- and check_en cannot see it either, since
% the English mirrors the German page for page and therefore inherits the same misalignment.
% Raised by two translation agents whose accounts of the two pages did not fit together.
Diese Sätze werden wesentlich verallgemeinert, wenn man eine lineare
Transformation darauf anwendet\footnote{\gesperrt{Laguerre}, Nouvelles
Annales de Math\'ematiques, 2.\,S\'er., Vol.\,19, 1880.}.

Dies wird leichter ausgeführt, wenn man die homogene Schreibweise anwendet,
also unter $f(x,y)$ eine ganze homogene Function $n^{\mathrm{ten}}$ Grades
versteht. Bildet man für einen beliebigen Werth $\xi:\eta$ die Polaren
dieser Function (\S.\,60):
\[
P_1(x,\xi) = \frac{1}{n}\{\xi f'(x) + \eta f'(y)\}
\]
\weq{3}{P_2(x,\xi) = \frac{1}{n(n-1)}\{\xi^2 f''(x,x) + 2\xi\eta f''(x,y) + \eta^2 f''(y,y)\}}
\[
P_3(x,\xi) = \frac{1}{n(n-1)(n-2)}\{\xi^3 f'''(x,x,x) + 3\xi^2\eta f'''(x,x,y) + 3\xi\eta^2 f'''(x,y,y) + \eta^3 f'''(y,y,y)\}
\]
\[
\cdots\cdots\cdots\cdots\cdots\cdots\cdots\cdots\cdots\cdots\cdots\cdots,
\]
so bleiben diese Functionen ungeändert, wenn man $\xi, \eta$ und $x, y$
gleichzeitig durch eine lineare Transformation umformt und für $f$ und
seine Ableitungen die entsprechenden transformirten Functionen setzt. In
Bezug auf die Realität der Wurzeln wird durch eine reelle lineare
Transformation nichts geändert. Nun lässt sich aber, wenn $\xi, \eta$ reell
sind, die reelle lineare Substitution so bestimmen, dass die transformirten
Werthe $\xi' = 1$, $\eta' = 0$ werden, und dadurch gehen, wenn man noch
$y = 1$ setzt, die Polaren in die Derivirten der Function $f(x)$ nach $x$
über (abgesehen von Zahlenfactoren).

\gesperrt{Daraus folgt, dass in den Sätzen XIII, XIV die Derivirten $f'(x), f''(x)
\dots$ durch die Polaren $P_1(x,\xi), P_2(x,\xi) \dots$ für ein beliebiges,
reelles $\xi, \eta$ ersetzt werden können.}

Wir machen hiervon die Anwendung auf die Gleichung, die man erhält, wenn
man die $(n-2)^{\mathrm{te}}$ Polare gleich Null setzt:
\weq{4}{x^2 f''(\xi,\xi) + 2xy f''(\xi,\eta) + y^2 f''(\eta,\eta) = 0.}
Hat $f(x,y) = 0$, wie wir voraussetzen, nur reelle Wurzeln, so muss auch die
quadratische Gleichung (4) reelle Wurzeln haben, d.\,h. es muss die
Hesse'sche Determinante
\weq{5}{f''(\xi,\xi) f''(\eta,\eta) - f''(\xi,\eta)^2}
für alle reellen $\xi, \eta$ negativ sein. Sie kann nur dann ver-
%CONT>>
